\section{Contexto y objetivo}
\label{sec:contexto}

El proyecto estima, para cada uno de los \textbf{170 municipios PDET}, cuántas
edificaciones hay y cuál es el área total de techos disponible para paneles
solares. Para eso compara dos fuentes abiertas de huellas de edificios y usa
una base de datos NoSQL con índices espaciales.

La Oficina de Investigación e Información de la UPME busca municipios
candidatos para proyectos piloto de granjas solares. A mayor
área de techos en un municipio, mayor potencial de captación. La lista oficial
de municipios viene en \campo{MunicipiosPDET.xlsx}: 170 municipios en 16
subregiones PDET y 19 departamentos~\cite{pdet}.

Esta primera entrega fija las bases del proyecto:

\begin{itemize}
  \item \textbf{Plan de implementación}: la metodología completa, desde la
        descarga de los datos hasta el informe final.
  \item \textbf{Diseño del esquema y modelo de datos}: el diagrama E-R de las
        colecciones y el diccionario de datos.
\end{itemize}

Las entregas siguientes construyen sobre este diseño: límites PDET (semana 2),
carga de edificios (semana 3) y análisis e informe final (semana 4).
